\documentclass[11pt]{letter}
\usepackage[margin=1in]{geometry}
\usepackage{amsmath,amssymb}
\usepackage{hyperref}

\signature{Bryan Ehrlich}
\address{Independent researcher\\bryan@ehrlich.dev}

\begin{document}

\begin{letter}{Editorial Board\\Journal of Mathematical Physics\\AIP Publishing}

\opening{Dear Editors,}

I am submitting the manuscript \emph{Faithful Self-Modeling Is Complex
Quantum Mechanics} for consideration as an Article in the Journal of
Mathematical Physics.

The paper proves a characterization theorem: among simple
finite-dimensional spectral order unit spaces, the self-modeling
systems (Definition~2.6) are exactly the complex matrix algebras
$M_n(\mathbb{C})^{\mathrm{sa}}$ equipped with the L\"uders sequential
product and conjugate-transpose involution.  The four conditions of
the definition---spectral state space, faithful tracking, minimal
composite, and irreducibility---unpack the concept of ``self-modeling
system'' without mentioning complex numbers, Hilbert spaces, or Born
probabilities; the complex field, $C^*$-involution, and full algebraic
structure of quantum theory emerge as consequences.

The main novel contribution is the construction of a sequential
product from the self-modeling cycle (Sections~3--4), where the key
step is a proof via the matrix Cauchy theorem that compatible
associativity forces scalar Peirce 1-space action and uniquely
determines the mixing function.  The downstream chain---EJA
classification, local tomography, type exclusion, $C^*$-algebra
promotion---relies on published theorems of van de Wetering, Barnum,
Wilce, Graydon, and Hanche-Olsen, and is machine-verified in Lean~4
with Mathlib (0 \texttt{sorry} declarations; code available at
\url{https://github.com/ehrlich-b/research/tree/jmp-submission-2026-03-28/lean}).

The paper contributes to the quantum reconstruction program and builds
on the sequential product framework of van de Wetering (JMP, 2019).
Suggested referees who work in this area include John van de Wetering,
Howard Barnum, Markus M\"uller, and Martin Pl\'avala.

The manuscript has not been submitted elsewhere.  AI tools were used
throughout: Claude (Anthropic) as an interactive research tool for
mathematical exploration, literature search, and manuscript drafting;
GPT-5.4 (OpenAI) for adversarial review of mathematical claims; and Lean~4 with Mathlib for machine verification of the
downstream derivation chain.  All definitions, theorems, and
scientific conclusions are my own.

\closing{Sincerely,}

\end{letter}
\end{document}
